\begin{table}[H]
\footnotesize
\vspace{-1em}
\caption{Sample multi-turn questions in MT-bench. }
\label{tab:mt_bench_sample}
\centering
\resizebox{0.9\columnwidth}{!}{
\begin{tabular}{@{}p{0.10\linewidth}|p{0.08\linewidth}p{0.75\linewidth}@{}}
\toprule
Category                  & \multicolumn{2}{c}{Sample Questions}                                                                                                                                                                                                                                                                           \\ \midrule
\multirow{2}{*}{Writing}  & 1st Turn & Compose an engaging travel blog post about a recent trip to Hawaii, highlighting cultural experiences and must-see attractions.                                                                                                                                                                        \\ \cmidrule(l){2-3} 
                          & 2nd Turn & Rewrite your previous response. Start every sentence with the letter A.                                                                                                                                                                                                                                \\ \midrule
\multirow{2}{*}{Math}     & 1st Turn & Given that $f(x) = 4x^3 - 9x - 14$, find the value of $f(2)$. \\ \cmidrule(l){2-3} 
                          & 2nd Turn & Find $x$ such that $f(x) = 0$. \\ \midrule
\multirow{2}{*}{Knowledge} & 1st Turn & Provide insights into the correlation between economic indicators such as GDP, inflation, and unemployment rates. Explain how fiscal and monetary policies ...                                                                                                                   \\ \cmidrule(l){2-3} 
                          & 2nd Turn & Now, explain them again like I’m five.                                                                                                                                                                                                                                                                 \\ \bottomrule
\end{tabular}
}
\vspace{-0.5em}
\end{table}

\section{MT-Bench and Chatbot Arena}
\label{sec:mt-bench}

\subsection{Motivation}
With the recent advances of LLMs, LLM-based assistants start to exhibit artificial general intelligence across diverse tasks, from writing and chatting to coding~\cite{bubeck2023sparks,openai2023gpt4,anil2023palm,srivastava2022beyond}. However, evaluating their broad capabilities also becomes more challenging. 
Despite the availability of numerous benchmarks for language models, they primarily focus on evaluating models on closed-ended questions with short responses. 
Given that these chat assistants can now precisely follow user instructions in multi-turn dialogues and answer open-ended questions in a zero-shot manner, current benchmarks are inadequate for assessing such capabilities. Existing benchmarks mostly fall into the following three categories.
\begin{itemize}
\item \textbf{Core-knowledge benchmarks}, including MMLU~\cite{hendrycks2020measuring}, HellaSwag~\cite{zellers2019hellaswag}, ARC~\cite{clark2018think}, WinoGrande~\cite{sakaguchi2021winogrande}, HumanEval~\cite{chen2021evaluating}, GSM-8K~\cite{cobbe2021training}, and AGIEval~\cite{zhong2023agieval}, evaluate the core capabilities of pre-trained LLMs using zero-shot and few-shot benchmark sets. They typically require LLMs to generate a short, specific answer to benchmark questions that can be automatically validated. %
\item \textbf{Instruction-following benchmarks}, such as  Flan~\cite{longpre2023flan,wei2021finetuned}, Self-instruct~\cite{selfinstruct}, NaturalInstructions~\cite{naturalinstructions}, Super-NaturalInstructions~\cite{supernaturalinstructions}, expand to slightly more open-ended questions and more diverse tasks and are used to evaluate LLMs after instruction fine-tuning.
\item \textbf{Conversational benchmarks}, like CoQA~\cite{reddy2019coqa}, MMDialog~\cite{feng2022mmdialog} and OpenAssistant~\cite{kopf2023openassistant}, are closest to our intended use cases. However, the diversity and complexity of their questions often fall short in challenging the capabilities of the latest chatbots.
\end{itemize}

While largely overlooked by existing LLM benchmarks, human preferences serve as a direct measure of a chatbot's utility in open-ended, multi-turn human-AI interactions. To bridge this gap, we introduce two novel benchmarks expressly tailored to assess human preferences. Simultaneously, these benchmarks are designed to distinguish the core capabilities of state-of-the-art models.


\subsection{MT-Bench}
We create MT-bench, a benchmark consisting of 80 high-quality multi-turn questions.
MT-bench is designed to test multi-turn conversation and instruction-following ability, covering common use cases and focusing on challenging questions to differentiate models.
We identify 8 common categories of user prompts to guide its construction: writing, roleplay, extraction, reasoning, math, coding, knowledge I (STEM), and knowledge II (humanities/social science).
For each category, we then manually designed 10 multi-turn questions.
\Cref{tab:mt_bench_sample} lists several sample questions.



\subsection{Chatbot Arena}
Our second approach is Chatbot Arena, a crowdsourcing benchmark platform featuring anonymous battles.
On this platform, users can interact with two anonymous models simultaneously, posing the same question to both. They vote for which model provides the preferred response, with the identities of the models disclosed post-voting. After running Chatbot Arena for one month, we have collected around 30K votes. Since the platform does not use pre-defined questions, it allows gathering a wide range of unrestricted use cases and votes in the wild, based on the diverse interests of users. A screenshot of the platform can be found at \Cref{subsec:chatbot-arena}.

